\documentclass[12pt,a4paper]{article}

\usepackage[utf8]{inputenc}
\usepackage[T1]{fontenc}
\usepackage[brazil]{babel}
\usepackage[a4paper,margin=2.5cm]{geometry}
\usepackage{booktabs}
\usepackage{array}
\usepackage{enumitem}
\usepackage{hyperref}
\hypersetup{hidelinks}

\setlength{\parindent}{1.5em}

\begin{document}

% ------------------------------------------------------------------ capa
\begin{titlepage}
\centering

{\large UNIVERSIDADE FEDERAL DO PARÁ (UFPA)}\\[0.3cm]
{\large INSTITUTO DE CIÊNCIAS EXATAS E NATURAIS (ICEN)}\\[0.3cm]
{\large CURSO DE CIÊNCIA DA COMPUTAÇÃO}\\[3.5cm]

{\LARGE\bfseries RELATÓRIO FINAL CENÁRIO DE REDE DE COMPUTADORES}\\[0.6cm]
{\large Segmentação por VLANs, Roteamento Inter-VLAN e DHCP Distribuído}\\[3.5cm]

{\large\bfseries Autores:}\\[0.3cm]
{\large Luiz Gabriel Santos Moreira - 202304940001}\\[0.2cm]
{\large Christian de Jesus da Costa Marinho - 202004940041}\\[3.5cm]

\vfill
{\large Belém - PA}\\[0.2cm]
{\large 2026}

\end{titlepage}

% ------------------------------------------------------------------ sumário
\tableofcontents
\newpage

% ------------------------------------------------------------------ 1
\section{INTRODUÇÃO}

Este relatório apresenta a análise e a documentação de uma infraestrutura de rede
corporativa departamentalizada, implementada e simulada na ferramenta Cisco Packet
Tracer (v9.0.1). O cenário exercita a segmentação de domínios de broadcast por meio
de VLANs (\textit{Virtual Local Area Networks}), o transporte de múltiplas VLANs
sobre enlaces de agregação com encapsulamento IEEE 802.1Q, o roteamento inter-VLAN
no modelo \textit{router-on-a-stick} e, na versão final aqui documentada, a
distribuição dinâmica de endereços IP por meio de servidores DHCP hospedados nos
próprios comutadores de acesso.

A relevância do cenário reside na reprodução de um problema recorrente em ambientes
empresariais: setores distintos precisam compartilhar a mesma infraestrutura física de
comutação sem compartilhar o mesmo domínio de broadcast. O isolamento de tráfego entre
departamentos é garantido pelas VLANs, enquanto a comunicação controlada entre eles é
provida pelo roteador central. A atribuição dinâmica de endereços, delegada aos
comutadores, elimina a dependência de um servidor DHCP externo dedicado e distribui
a carga entre os dois equipamentos de acesso.

\subsection{Contextualização e Motivação do Cenário Corporativo}

A arquitetura foi idealizada para atender a uma organização de médio porte dividida
em quatro setores funcionais, cada um mapeado em uma VLAN e em uma sub-rede IP
próprias, acrescida de uma VLAN de gerência/suporte:

\begin{itemize}[leftmargin=2em]
  \item \textbf{VLAN 10 --- VENDAS:} setor comercial. Todos os terminais utilizam
        endereçamento dinâmico via DHCP, refletindo a rotatividade característica
        de equipes de vendas e de dispositivos compartilhados como impressoras.

  \item \textbf{VLAN 20 --- MARKETING:} setor de comunicação e campanhas, com
        endereçamento estático atribuído manualmente às estações.

  \item \textbf{VLAN 30 --- ADM:} setor administrativo. Por tratar dados sensíveis
        de gestão, exige isolamento em relação aos demais domínios de broadcast.

  \item \textbf{VLAN 40 --- TIC:} setor de Tecnologia da Informação e Comunicação,
        responsável pela operação da própria infraestrutura de rede.

  \item \textbf{VLAN 1 --- Gerência/Suporte:} reservada ao tráfego de gerência dos
        ativos de camada 2 por meio de interfaces virtuais de comutação (SVI) e ao
        acesso de laptops de suporte técnico que precisam alcançar e configurar
        qualquer equipamento da rede.
\end{itemize}

% ------------------------------------------------------------------ 2
\section{ARQUITETURA DA SOLUÇÃO}

A arquitetura física adota um modelo hierárquico de dois níveis. A camada de acesso é
formada por dois comutadores Cisco Catalyst 2960-24TT (\textbf{SW0} e \textbf{SW1}),
que concentram as estações de trabalho, os notebooks e as impressoras dos quatro
setores. A camada de distribuição e roteamento é exercida por um único roteador Cisco
(\textbf{Router0}), conectado à malha de comutação pelo enlace de \textit{uplink}
proveniente de \textbf{SW1}. Um servidor dedicado (\textbf{Server0}) está conectado
diretamente à interface \texttt{GigabitEthernet0/0} do roteador, reservado para
serviços futuros.

Cada VLAN de produção está distribuída simetricamente entre os dois comutadores: as
mesmas faixas de portas físicas foram atribuídas às mesmas VLANs em ambos os
equipamentos, e o enlace de agregação entre eles mantém a continuidade de cada domínio
de broadcast.

\subsection{Mapeamento de Sub-redes e Segmentação Lógica}

\begin{sloppypar}
O esquema de endereçamento adota blocos privados de Classe C com máscara \texttt{/24}
(\texttt{255.255.255.0}), alinhando o terceiro octeto ao identificador numérico da
VLAN correspondente. As sub-redes são:

\begin{itemize}[leftmargin=2em]
  \item \textbf{VLAN 1 --- Gerência/Suporte:} \texttt{192.168.1.0/24}. SVIs dos
        comutadores em \texttt{192.168.1.1} (SW0) e \texttt{192.168.1.2} (SW1).
        Gateway em \texttt{192.168.1.254} (subinterface do roteador). Laptops de
        suporte recebem IPs via DHCP.

  \item \textbf{VLAN 10 --- VENDAS:} \texttt{192.168.10.0/24}. Gateway em
        \texttt{192.168.10.1}. Todos os terminais utilizam DHCP; SW0 serve a
        faixa \texttt{.11}--\texttt{.100} e SW1 a faixa \texttt{.101}--\texttt{.200}.

  \item \textbf{VLAN 20 --- MARKETING:} \texttt{192.168.20.0/24}. Gateway em
        \texttt{192.168.20.254}. Endereçamento estático: SW0 (\texttt{.7}--\texttt{.9}),
        SW1 (\texttt{.107}--\texttt{.109}). Pools DHCP provisionados para
        expansão futura.

  \item \textbf{VLAN 30 --- ADM:} \texttt{192.168.30.0/24}. Gateway em
        \texttt{192.168.30.254}. Endereçamento estático: SW0 (\texttt{.13}--\texttt{.15}),
        SW1 (\texttt{.113}--\texttt{.115}). Pools DHCP provisionados.

  \item \textbf{VLAN 40 --- TIC:} \texttt{192.168.40.0/24}. Gateway em
        \texttt{192.168.40.254}. Endereçamento estático: SW0 (\texttt{.19}--\texttt{.21}),
        SW1 (\texttt{.119}--\texttt{.121}). Pools DHCP provisionados.
\end{itemize}
\end{sloppypar}

\subsection{Topologia de Terminais por Setor e Comutador}

A distribuição simétrica de dispositivos entre SW0 e SW1 é detalhada a seguir:

\begin{itemize}[leftmargin=2em]
  \item \textbf{VLAN 10 --- VENDAS:} PC3, PC4, PC5 e Printer0 em SW0
        (Fa0/1--Fa0/4); Laptop2(2), Laptop3(2), Laptop4(2) e Printer1 em SW1
        (Fa0/1--Fa0/4). Endereçamento exclusivamente via DHCP.

  \item \textbf{VLAN 20 --- MARKETING:} PC9, PC10, PC11 em SW0 (Fa0/7--Fa0/9);
        Laptop2(3), Laptop3(3), Laptop4(3) em SW1 (Fa0/7--Fa0/9).

  \item \textbf{VLAN 30 --- ADM:} PC8, PC7, PC6 em SW0 (Fa0/13--Fa0/15);
        Laptop2(1), Laptop3(1), Laptop4(1) em SW1 (Fa0/13--Fa0/15).

  \item \textbf{VLAN 40 --- TIC:} PC0, PC1, PC2 em SW0 (Fa0/19--Fa0/21);
        Laptop4, Laptop3, Laptop2 em SW1 (Fa0/19--Fa0/21).

  \item \textbf{VLAN 1 --- Suporte:} Laptop0 em SW0 (Fa0/24) e Laptop5 em SW1
        (Fa0/24). Endereçamento via DHCP; SW0 serve \texttt{.10}--\texttt{.50} e
        SW1 serve \texttt{.51}--\texttt{.100}.
\end{itemize}

\subsection{Papel dos Ativos de Rede}

\begin{itemize}[leftmargin=2em]
  \item Os dois comutadores \textbf{2960-24TT} operam na \textbf{Camada 2}, segregando
        quadros Ethernet em domínios de broadcast distintos e transportando-os marcados
        pelos enlaces de agregação. Adicionalmente, exercem o papel de
        \textbf{servidores DHCP} para todas as VLANs, servindo faixas de endereços
        complementares e sem sobreposição.

  \item As SVIs \texttt{Vlan1} de ambos os comutadores
        (\texttt{192.168.1.1} em SW0 e \texttt{192.168.1.2} em SW1) exercem a
        função de \textbf{gerência} do equipamento e servem como endereço de destino
        para o relay DHCP das VLANs de produção.

  \item O \textbf{roteador} (Router0) é o único ativo de \textbf{Camada 3}. Opera
        no modelo \textit{router-on-a-stick}: cinco subinterfaces lógicas sobre
        \texttt{GigabitEthernet2/0} exercem o papel de \textit{default gateway}
        de cada VLAN e encaminham pedidos DHCP aos servidores nos comutadores
        por meio do mecanismo \texttt{ip helper-address}.

  \item O protocolo \textbf{Spanning Tree} opera em modo \texttt{pvst} nos dois
        comutadores, com uma instância independente por VLAN, prevenindo laços de
        camada 2 no enlace de agregação.
\end{itemize}

% ------------------------------------------------------------------ 3
\section{IMPLEMENTAÇÃO}

\subsection{Base de VLANs e Portas de Acesso}

As quatro VLANs de produção foram criadas de forma idêntica e independente na base de
dados local de cada comutador, com os nomes \texttt{VENDAS}, \texttt{MARKETING},
\texttt{ADM} e \texttt{TIC}. Ambos os equipamentos operam com VTP em modo
\textit{server}. Todas as portas de acesso receberam o comando
\texttt{switchport mode access} de forma explícita, eliminando a negociação DTP.
As portas não utilizadas foram desativadas com \texttt{shutdown}.

\subsection{DHCP Distribuído nos Comutadores}

A arquitetura de DHCP adota dois servidores redundantes e complementares, hospedados
nos próprios comutadores de acesso. Cada switch mantém pools para \textbf{todas as
VLANs} (1, 10, 20, 30 e 40), mas serve faixas de endereços distintas, evitando
conflitos mesmo quando ambos respondem ao mesmo pedido:

\begin{itemize}[leftmargin=2em]
  \item \textbf{SW0} distribui IPs na faixa \texttt{.10}--\texttt{.50} para VLAN 1
        e na faixa \texttt{.11}--\texttt{.100} para as VLANs 10 a 40.
  \item \textbf{SW1} distribui IPs na faixa \texttt{.51}--\texttt{.100} para VLAN 1
        e na faixa \texttt{.101}--\texttt{.200} para as VLANs 10 a 40.
\end{itemize}

As exclusões são implementadas com \texttt{ip dhcp excluded-address}, garantindo que
cada switch ofereça apenas sua metade da faixa disponível. Os endereços de gateway,
das SVIs e de infraestrutura estão excluídos em ambos os servidores.

\subsection{Relay DHCP no Roteador}

Para que pedidos DHCP originados nas VLANs de produção (10 a 40) cheguem aos
servidores hospedados nas SVIs dos comutadores (\texttt{192.168.1.1} e
\texttt{192.168.1.2}), cada subinterface do roteador recebe dois comandos
\texttt{ip helper-address}:

\begin{itemize}[leftmargin=2em]
  \item \texttt{ip helper-address 192.168.1.1} --- encaminha para SW0
  \item \texttt{ip helper-address 192.168.1.2} --- encaminha para SW1
\end{itemize}

\begin{sloppypar}
Para que o roteador consiga encaminhar esses pacotes unicast aos endereços das SVIs,
foi necessário criar a subinterface \texttt{GigabitEthernet2/0.1} com o endereço
\texttt{192.168.1.254/24}, inserindo a rede \texttt{192.168.1.0/24} na tabela de
roteamento. Para a VLAN~1, o DHCP funciona sem relay, pois os próprios comutadores
estão nessa VLAN e respondem diretamente aos broadcasts.
\end{sloppypar}

\subsection{Subinterfaces do Roteador}

O roteador possui cinco subinterfaces configuradas sobre \texttt{GigabitEthernet2/0}:

\begin{itemize}[leftmargin=2em]
  \item \texttt{Gi2/0.1}: \texttt{encapsulation dot1Q 1 native},
        \texttt{192.168.1.254/24} --- gateway e ponto de relay da VLAN 1.
  \item \texttt{Gi2/0.10}: \texttt{encapsulation dot1Q 10},
        \texttt{192.168.10.1/24} --- gateway VENDAS + relay DHCP.
  \item \texttt{Gi2/0.20}: \texttt{encapsulation dot1Q 20},
        \texttt{192.168.20.254/24} --- gateway MARKETING + relay DHCP.
  \item \texttt{Gi2/0.30}: \texttt{encapsulation dot1Q 30},
        \texttt{192.168.30.254/24} --- gateway ADM + relay DHCP.
  \item \texttt{Gi2/0.40}: \texttt{encapsulation dot1Q 40},
        \texttt{192.168.40.254/24} --- gateway TIC + relay DHCP.
\end{itemize}

\subsection{Enlace de Agregação 802.1Q}

Dois enlaces de agregação transportam as VLANs 1, 10, 20, 30 e 40:

\begin{itemize}[leftmargin=2em]
  \item \textbf{Enlace entre comutadores:} \texttt{GigabitEthernet0/1} de SW0
        e SW1 (cabo cruzado), garantindo a continuidade dos domínios de broadcast
        entre as duas caixas de acesso.
  \item \textbf{Enlace de \textit{uplink}:} \texttt{GigabitEthernet0/2} de SW1
        até \texttt{GigabitEthernet2/0} do Router0. Todo o tráfego inter-VLAN e
        o relay DHCP transitam por este enlace.
\end{itemize}

A porta \texttt{Gi0/2} de SW0 permanece em \texttt{shutdown}.

\subsection{Laptops de Suporte (VLAN 1)}

\begin{sloppypar}
Os laptops Laptop0 (SW0, Fa0/24) e Laptop5 (SW1, Fa0/24)
integram a VLAN~1 de gerência. Com endereços obtidos via DHCP na sub-rede
\texttt{192.168.1.0/24} e gateway \texttt{192.168.1.254}, esses dispositivos
de suporte técnico alcançam as SVIs de ambos os comutadores
(\texttt{.1} e \texttt{.2}) e as subinterfaces do roteador, podendo configurar
qualquer equipamento da infraestrutura por acesso remoto.
\end{sloppypar}

% ------------------------------------------------------------------ 4
\section{ANÁLISE DO CENÁRIO}

O cenário implementado é funcional e plenamente operacional para os objetivos
propostos. A seguir registram-se as observações analíticas relevantes.

\begin{enumerate}[leftmargin=2em]
  \item \textbf{Roteamento inter-VLAN completo.} Todas as cinco subinterfaces
        do roteador estão configuradas e ativas. Qualquer terminal com IP válido
        pode alcançar terminais de qualquer outra VLAN, bastando que o \textit{gateway}
        esteja corretamente configurado.

  \item \textbf{DHCP redundante e sem sobreposição.} A divisão de faixas entre SW0
        e SW1 garante que dois servidores DHCP possam responder ao mesmo pedido sem
        conflito de endereços. O host aceita a primeira oferta que chega,
        independentemente de qual switch respondeu.

  \item \textbf{Relay DHCP dependente da subinterface Gi2/0.1.} O correto
        funcionamento do relay exige que a rede \texttt{192.168.1.0/24} esteja
        presente na tabela de roteamento do Router0. A ausência desta subinterface
        impede o encaminhamento unicast dos pedidos DHCP às SVIs dos comutadores.

  \item \textbf{Servidor sem endereçamento.} O Server0, conectado à
        \texttt{GigabitEthernet0/0} do roteador, não possui endereço IP
        e a interface permanece sem endereçamento (\texttt{no ip address}).
        O segmento de servidores está inoperante.

  \item \textbf{Endurecimento do acesso administrativo pendente.} As linhas
        \texttt{vty} dos três ativos exigem autenticação (\texttt{login}), mas
        nenhuma senha foi definida, inviabilizando o acesso remoto por Telnet/SSH.
        Não há \texttt{enable secret} configurado. A diretiva
        \texttt{ip default-gateway 192.168.1.254} no SW0 está correta após a
        adição da subinterface Gi2/0.1 no roteador.

  \item \textbf{Coexistência de dois servidores VTP.} Ambos os comutadores operam
        em modo \textit{server} com revisões independentes. Recomenda-se adotar
        o modo \textit{transparent} ou definir um único domínio VTP para evitar
        sobreposição involuntária das bases de VLANs.

  \item \textbf{Ausência de caminho redundante para o roteador.} A porta
        \texttt{Gi0/2} de SW0 encontra-se desativada. Todo o tráfego inter-VLAN
        passa obrigatoriamente pelo enlace SW1~$\rightarrow$~Router0, configurando
        ponto único de falha. A ativação de \texttt{Gi0/2} em SW0 e a adição de
        uma segunda subinterface no roteador permitiriam redundância protegida
        pelo Spanning Tree.
\end{enumerate}

% ------------------------------------------------------------------ 5
\section{RESUMO DO CENÁRIO}

\begin{table}[htbp]
\centering
\caption{Plano de Endereçamento Global por VLAN}
\label{tab:enderecamento}
\small
\begin{tabular}{@{}clllll@{}}
\toprule
\textbf{VLAN} & \textbf{Nome} & \textbf{Rede} & \textbf{Gateway} &
\textbf{Faixa SW0} & \textbf{Faixa SW1} \\
\midrule
1  & Gerência  & 192.168.1.0/24  & 192.168.1.254   & .10–.50   & .51–.100  \\
10 & VENDAS    & 192.168.10.0/24 & 192.168.10.1    & .11–.100  & .101–.200 \\
20 & MARKETING & 192.168.20.0/24 & 192.168.20.254  & .11–.100  & .101–.200 \\
30 & ADM       & 192.168.30.0/24 & 192.168.30.254  & .11–.100  & .101–.200 \\
40 & TIC       & 192.168.40.0/24 & 192.168.40.254  & .11–.100  & .101–.200 \\
\bottomrule
\end{tabular}
\end{table}

\begin{table}[htbp]
\centering
\caption{Alocação de Portas nos Comutadores de Acesso (padrão idêntico em ambos)}
\label{tab:portas}
\small
\begin{tabular}{@{}p{4.0cm}p{2.4cm}p{2.4cm}p{3.8cm}@{}}
\toprule
\textbf{Portas} & \textbf{VLAN} & \textbf{Modo} & \textbf{Observação} \\
\midrule
Fa0/1 -- Fa0/4              & 10 (VENDAS)    & Access explícito  & PCs e impressoras \\
Fa0/7 -- Fa0/9              & 20 (MARKETING) & Access explícito  & PCs e notebooks \\
Fa0/13 -- Fa0/15            & 30 (ADM)       & Access explícito  & PCs e notebooks \\
Fa0/19 -- Fa0/21            & 40 (TIC)       & Access explícito  & PCs e notebooks \\
Fa0/24                      & 1 (Gerência)   & Access padrão     & Laptop de suporte \\
Fa0/5, Fa0/6, Fa0/10--12   & ---            & \texttt{shutdown} & Portas de reserva \\
Fa0/16--18, Fa0/22--23      & ---            & \texttt{shutdown} & Portas de reserva \\
Gi0/1                       & 1,10,20,30,40  & \textit{Trunk}    & Enlace inter-switches \\
Gi0/2 (SW1)                 & 1,10,20,30,40  & \textit{Trunk}    & \textit{Uplink} ao roteador \\
Gi0/2 (SW0)                 & ---            & \texttt{shutdown} & Porta sem uso \\
\bottomrule
\end{tabular}
\end{table}

\begin{table}[htbp]
\centering
\caption{Subinterfaces do Roteador (Router0)}
\label{tab:subinterfaces}
\small
\begin{tabular}{@{}p{3.8cm}cllp{3.2cm}@{}}
\toprule
\textbf{Subinterface} & \textbf{VLAN} & \textbf{Endereço IP} &
\textbf{Máscara} & \textbf{Função} \\
\midrule
GigabitEthernet2/0.1  & 1  & 192.168.1.254   & 255.255.255.0 & Gateway + relay VLAN 1 \\
GigabitEthernet2/0.10 & 10 & 192.168.10.1    & 255.255.255.0 & Gateway + relay VENDAS \\
GigabitEthernet2/0.20 & 20 & 192.168.20.254  & 255.255.255.0 & Gateway + relay MARKETING \\
GigabitEthernet2/0.30 & 30 & 192.168.30.254  & 255.255.255.0 & Gateway + relay ADM \\
GigabitEthernet2/0.40 & 40 & 192.168.40.254  & 255.255.255.0 & Gateway + relay TIC \\
\bottomrule
\end{tabular}
\end{table}

\clearpage

% ------------------------------------------------------------------ 6
\section{REFERÊNCIAS}

\begin{itemize}[leftmargin=2em]
  \item CISCO SYSTEMS. \textit{Cisco Packet Tracer}. Disponível em:
        \url{https://www.netacad.com/courses/packet-tracer}. Acesso em: setembro. 2026.

  \item IEEE. \textit{IEEE Standard for Local and Metropolitan Area Networks ---
        Bridges and Bridged Networks (IEEE 802.1Q)}. Nova York: IEEE, 2018.

  \item TANENBAUM, Andrew S.; WETHERALL, David J. \textit{Redes de Computadores}.
        5. ed. São Paulo: Pearson, 2011.

  \item KUROSE, James F.; ROSS, Keith W. \textit{Redes de Computadores e a Internet:
        Uma Abordagem Top-Down}. 7. ed. São Paulo: Pearson, 2013.

  \item CISCO NETWORKING ACADEMY. \textit{CCNA: Switching, Routing, and Wireless
        Essentials}. Cisco Press, 2020.
\end{itemize}

\end{document}
